% ===================================================================
%  TABELLA N-RO 13 — Le hotel. — Le restaurante. — Le caffe.
%  Traduction 2026 d'apres texto/io, controlee sur texto/fr et
%  texto/es. Consigne : voir texto/ia/10-tabella-01.tex.
%
%  Deux notes, marquees toutes deux « (*) », et deux appels : celui du
%  bloc 01-2 (la chambre) et celui du bloc 09-1 (le menu). L'ordre les
%  departage.
% ===================================================================

%%K t13-noto-1 noto
\VUnotes{143.2mm}{%
\textit{\VUgras{(*) Quanto al detalios del camera, vide le tabella
N\textsuperscript{o}. 6.}}%
}

%%K t13-apar-1 apar
\VUsaut{3.0mm}
\VUtitre{15.0pt}{60}{TABELLA N\textsuperscript{o} 13}
\VUsaut{1.6mm}
\VUpk{10.2pt}{40}{Le hotel. --- Le restaurante.}
\VUsaut{2.0mm}
\VUpk{10.2pt}{40}{Le caffe.}
\VUsaut{3.4mm}

%%K t13-01-1 p
1. --- Habente arrivate heri vespere in Royan, io me alogiava al « \textit{Hotel
del Oceano} » que me habeva essite recommendate per un amico.

%%K t13-01-2 p
On me forniva un belle \VUgras{camera} \textsuperscript{(2)}, sur le secunde
\VUgras{etage} \textsuperscript{(1)} ubi io dormiva multo ben (*).

%%K t13-02-1 p
2. --- Iste matino, exiente de mi camera, in le qual le \VUgras{servitrice}
\textsuperscript{(3)} habeva apportate le parve dejuner, io entrava le
\VUgras{fumatorio} \textsuperscript{(5)} que es anque usate como \VUgras{sala de
lectura} \textsuperscript{(4)} pro leger le \VUgras{jornales}
\textsuperscript{(6)}, fumante un \VUgras{cigarretta} \textsuperscript{(8)}.
Post leger, io descendeva per le \VUgras{ascensor} \textsuperscript{(9)} e iva
promenar tra le citate.

%%K t13-03-1 p
3. --- Proque io habeva oblidate questionar le \VUgras{empleato}
\textsuperscript{(12)} del hotel super le hora quando on da le repastos sur le
\VUgras{tabula commun} \textsuperscript{(11)}, io retornava tosto e profitava lo
pro ir banar me in le \VUgras{sala de banio} \textsuperscript{(10)} del hotel.
Postea io iva de novo al \VUgras{cassa} \textsuperscript{(15)} pro telephonar a
un de mi amicos. Ma le \VUgras{telephono} \textsuperscript{(13)} esseva occupate;
vidente mi embarasso le \VUgras{cassera} \textsuperscript{(14)}, qui registrava
un \VUgras{factura} \textsuperscript{(17)} sur le \VUgras{registro del viageros}
\textsuperscript{(16)} rogava le \VUgras{portero} \textsuperscript{(18)} inviar
le \VUgras{groom} \textsuperscript{(19)} facer mi commission \VUgras{in
bicycletta} \textsuperscript{(20)}.

%%K t13-04-1 p
4. --- Io la gratiava e exiva pro dar un gratification al \VUgras{portator}
\textsuperscript{(24)} (le portator de fardellos) qui apportava del station sur
su \VUgras{carretta a mano} \textsuperscript{(25)} mi \VUgras{cassa de cappellos}
\textsuperscript{(26)}, \VUgras{mi valise} \textsuperscript{(27)} e \VUgras{mi
sacco de viage} \textsuperscript{(28)}. Le \VUgras{portator de litteras}
\textsuperscript{(21)} justo arrivate, me dava un \VUgras{littera}
\textsuperscript{(22)}.

%%K t13-05-1 p
5. --- Io restava ancora un pauco de tempore sur le limine del porta, presso le
\VUgras{cassa a litteras} \textsuperscript{(23)}, reguardante le circulation sur
le strata. Le \VUgras{conductor} \textsuperscript{(30)} del \VUgras{omnibus}
\textsuperscript{(29)} cargava sur su vehiculo le \VUgras{cofro}
\textsuperscript{(31)} de un \VUgras{viagero} \textsuperscript{(32)} que le
\VUgras{patron del hotel} \textsuperscript{(33)} adeuava. Un juvene
\VUgras{telegraphista} \textsuperscript{(35)} sin hasta apportava un
\VUgras{telegramma} \textsuperscript{(34)} pro un cliente del hotel; un
\VUgras{turista} \textsuperscript{(36)} con su \VUgras{apparato photographic}
\textsuperscript{(38)} in bandoliera, consultava su \VUgras{guida}
\textsuperscript{(37)}.

%%K t13-tit-1 sub
\VUsaut{4.0mm}
\VUpk{10.2pt}{40}{Hora del dejuner.}
\VUsaut{3.0mm}

%%K t13-06-1 p
6. --- Ante le grillia, un insouciante \VUgras{commissionario}
\textsuperscript{(39)} dormettava inter su \VUgras{croco de portator}
\textsuperscript{(40)} e su \VUgras{cassa de cera} \textsuperscript{(41)},
durante que un juvene \VUgras{lavandera} \textsuperscript{(42)} qui portava un
\VUgras{corbe} \textsuperscript{(43)} de linage, rideva del solemne mina del
\VUgras{maestro de sala} \textsuperscript{(44)} qui stava presso le porta.

%%K t13-07-1 p
7. --- Vidente le \VUgras{cocinero} \textsuperscript{(45)}, qui exiva de su
\VUgras{cocina} \textsuperscript{(47)} situate al \VUgras{sub-solo}
\textsuperscript{(48)} pro inviar un \VUgras{lavaplattos} \textsuperscript{(46)}
facer un commission, io memorava que le hora del dejuner es proxime e io
entrava le restaurante, transversante le corte in le qual un \VUgras{automobilista}
\textsuperscript{(50)} garava su \VUgras{automobile} \textsuperscript{(49)},
durante que un \VUgras{machinista} \textsuperscript{(52)} reparava, secundo le
indicationes del \VUgras{cyclista} \textsuperscript{(53)}, un \VUgras{tricyclo a
petroleo} \textsuperscript{(51)} que justo habeva un accidente.

%%K t13-08-1 p
8. --- Io questionava un juvene \VUgras{groom} \textsuperscript{(54)} qui portava
\VUgras{vitros de bira} \textsuperscript{(55)}, si le tabula commun es sub le
veranda, e pro su responsa affirmative, io occupava un loco a un del tabulas e
faceva dar me un excellente repasto.

%%K t13-noto-2 noto
\VUnotes{143.2mm}{%
\textit{\VUgras{(*) Con le adjuta del lista de plattos introducite in le
vocabulario, le professor potera facilemente diversificar le menu sequente.}}%
}

%%K t13-tit-2 sub
\VUsaut{4.0mm}
\VUpk{10.2pt}{40}{Excellente dejuner.}
\VUsaut{3.0mm}

%%K t13-09-1 p
9. --- Le garson me apportava le menu (*) del dejuner e le lista de vinos. Io
eligeva e mandava como \VUgras{hors-d'oeuvre gambettas} con \VUgras{butyro}, e,
como entrata, \VUgras{tripas al modo de Caen}. Io non mangiava \VUgras{pisce},
ma demandava un portion de \VUgras{gamba} de agno, e, como \VUgras{platto de
legumines, spinacia} con crema. Postea, proque nulle del \VUgras{intermissos}
me placeva, io faceva apportar diverse dessertos, \VUgras{caseo},
\VUgras{fructos} e \VUgras{biscuits}. Io in plus adjungeva un excellente
Saint-Émilion-\VUgras{vino}, e multo contente de mi repasto io lassava le tabula
post haber donate al \VUgras{garson} \textsuperscript{(57)} un belle
\VUgras{gratification} \textsuperscript{(59)}.

%%K t13-10-1 p
10. --- Vidente me preste a partir, le \VUgras{gerente} \textsuperscript{(60)} me
proponeva ir al balcon pro admirar le splendide panorama del \VUgras{Oceano}
\textsuperscript{(62)}. Lontano on comenciava a vider \VUgras{parve barcas} de
pisca \textsuperscript{(63)}, \VUgras{naves a vela} \textsuperscript{(64)} e un
enorme \VUgras{paquebot} \textsuperscript{(65)} del qual le nigre silhouette se
distacca sur le blanc \VUgras{nubes} \textsuperscript{(67)} del \VUgras{horizonte}
\textsuperscript{(66)}. Un legier brisa faceva undular le \VUgras{standardo}
\textsuperscript{(70)} fixate super le \VUgras{insignia} \textsuperscript{(69)}
del \VUgras{hall} \textsuperscript{(68)}.

%%K t13-tit-3 sub
\VUsaut{3.0mm}
\VUpk{10.2pt}{40}{Divertimentos.}
\VUsaut{3.0mm}

%%K t13-11-1 p
11. --- Le spectaculo esseva tanto interessante, que io decideva ascender deman
sur le terrassa del caffe portante un \VUgras{binoculo} \textsuperscript{(71)} o
un \VUgras{telescopio} \textsuperscript{(72)}.

%%K t13-12-1 p
12. --- Lassante le balcon, io iva al caffe del hotel pro glutir un
\VUgras{bibita} \textsuperscript{(73)}, jocante un \VUgras{partita de billiard}
\textsuperscript{(75)}. Sin haltar io transversava le salon del caffe in le qual
multe consumitores jocava al \VUgras{chacos} \textsuperscript{(78)}, al
\VUgras{damas} \textsuperscript{(79)}, al \VUgras{dominos}
\textsuperscript{(80)}, al \VUgras{trictrac} \textsuperscript{(81)} o al
\VUgras{cartas} \textsuperscript{(82)}, e io ascendeva directemente al
\VUgras{sala de billiard} \textsuperscript{(74)}.

%%K t13-13-1 p
13. --- Quando le partita esseva finite, io descendeva al terreno e demandava al
garson un \VUgras{inveloppe} \textsuperscript{(84)}, \VUgras{papiro}
\textsuperscript{(83)}, un \VUgras{timbro postal} \textsuperscript{(85)}, un
\VUgras{penna} \textsuperscript{(87)} e \VUgras{incausto} \textsuperscript{(86)}
pro scriber un littera.

%%K t13-13-2 p
Postea io appellava un \VUgras{cochero} \textsuperscript{(92)} del qual le
\VUgras{vehiculo} \textsuperscript{(88)} habeva haltate ante le caffe. Io le
questionava super le precio secundo le horas o un cursa, e, quando nos
convenieva super le precio, ille me aperiva le \VUgras{porta del vehiculo}
\textsuperscript{(89)} e me installava. Ille reascendeva sur le \VUgras{sede}
\textsuperscript{(91)}, rapidemente faceva partir le cavallos con un
\VUgras{parve flagello} \textsuperscript{(93)} e nos partiva pro un incantatori
promenada.
\VUsaut{6.0mm}
\VUfilet{22mm}
